\section{Discussion and Conclusion}

Three broad lessons emerge from this study.

\textbf{Exposure bias is a first-order concern for autoregressive
cellular-automaton models, and its fix must be chosen carefully.} Teacher
forcing alone (V5) produces excellent single-step metrics that do not
predict rollout stability. The most direct-seeming fix for multi-step
training --- a straight-through estimator (V6) --- fails due to a
systematic directional gradient bias rather than an optimization-stability
issue, and is not rescued by gradient clipping or a lower learning rate.
Scheduled sampling (V7), which avoids passing gradients across unrolled
steps entirely, resolves the problem cleanly. This distinction --- between
instabilities caused by gradient \emph{magnitude} and those caused by
gradient \emph{direction} --- recurred later in the polynomial-activation
negative result (Section~\ref{sec:v11}), where we again observed a
divergence with a similar signature (stable training loss, exploding
validation loss) under a different mechanism (an unbounded activation
interacting with autoregressive feedback rather than an STE gradient path).

\textbf{Aggregate validation metrics can hide severe, structured
regressions.} The V8-to-V9 comparison (Section~\ref{sec:training-evolution})
is the clearest demonstration: V9's validation loss and precision/recall
are statistically indistinguishable from V8's, yet its per-grid error count
at high density is an order of magnitude worse. This was only detectable
via a targeted, density-stratified failure search --- a general reminder
that i.i.d.\ validation splits can fail to surface regressions concentrated
in identifiable input subpopulations, and that practitioners should audit
per-subgroup behavior rather than relying on a single scalar validation
metric, especially before deciding that a change (like scaling up training
data) has helped.

\textbf{Once a model is capacity-limited, only capacity fixes the problem.}
The clearest evidence for this is negative: V9 attempted to resolve V8's
residual high-density errors purely with more data (of a kind the model
already saw extensively) and made things worse. V10 held the data
distribution fixed and scaled the model 5.15$\times$, and eliminated all
2,500 evaluated per-grid errors. This is consistent with, though not
identical to, the picture in Ahmed \& Davis \cite{ahmed2026polynomial}: they
show that a poorly-matched inductive bias (ReLU) can make a
network fail to learn a rule it is representationally capable of encoding
with very few parameters, while we show that once a specific
\emph{well-matched} architecture (GELU CNN-Transformer with scheduled
sampling) has saturated its capacity at a given width, adding width is what
resolves the remaining gap --- not more data of the same kind. Our
attempt to combine both ideas (Section~\ref{sec:v11}) was unsuccessful in
its first iteration, underscoring that these two levers do not necessarily
compose for free, particularly once autoregressive training is involved.

\textbf{The embedding analysis (Section~\ref{sec:embedding-analysis})}
provides a mechanistic complement to the accuracy results: it shows
directly, via patch-level cosine similarity and PCA projections, that the
model's learned absolute positional embedding is responsible for most of
its sensitivity to geometric transformations on dense grids, while sparse
configurations are nearly embedding-invariant simply because most patches
are empty on both sides of any transformation. This is consistent with the
architectural choice (Section~\ref{sec:architecture}) to augment with the
dihedral group $D_4$ but not with translation, since translation
augmentation would directly conflict with a learnable absolute positional
embedding.

\subsection*{Limitations and future work}

This report reflects a single training run per configuration for the
CNN-Transformer experiments (V5--V11), aside from the deliberately
multi-seed reproduction in Section~\ref{sec:reproduction}; conclusions
about V6/V9/V11 instability would be strengthened by repeating those runs
across seeds, as we did for the minimal-CNN verification. The V11 negative
result (Section~\ref{sec:v11}) was not root-caused beyond the hypothesis
offered; a natural next step is to test the polynomial activation in a
single-step-only (teacher-forcing) training regime at V8/V10 model scale,
isolating whether the instability is intrinsic to the activation at this
model size or specific to its interaction with scheduled sampling. The
per-grid failure-search protocol used throughout
(Section~\ref{sec:training-evolution}) covers five densities and a small
set of named patterns; a more exhaustive audit (finer density grid,
targeted adversarial search) might reveal further tail-case regressions
not captured here, including for V10 despite its zero-error result on the
2,500 grids tested.
